\documentclass[11pt]{article}

\usepackage{amsmath,amssymb,amstext,amsthm}
\usepackage{geometry}
\usepackage{fancyhdr}
\usepackage[pdftex]{graphicx}
\usepackage{xcolor}

\geometry{
  verbose,
  dvips,
  width=430pt, marginparsep=0pt, marginparwidth=0pt,
  top=90pt, headheight=12pt, headsep=20pt, footskip=30pt, bottom=80pt}

\setlength{\parskip}{\medskipamount}
\setlength{\parindent}{0pt}

\linespread{1.05}

\newtheorem{theorem}{Theorem}
\newtheorem{lemma}[theorem]{Lemma}
\newtheorem{proposition}[theorem]{Proposition}
\newtheorem{corollary}[theorem]{Corollary}
\newtheorem{conjecture}[theorem]{Conjecture}
\newtheorem{obs}{Observation}
\newtheorem{clm}{Claim}
\newtheorem{ques}{Question}
\newtheorem{problem}{Problem}

\newtheorem{ex}{Example}


\newcommand{\spc}{\hspace{0.08in}}
\newcommand{\vs}{\vspace*{.1in}}
\newcommand{\E}{\mathbb{E}}
\newcommand{\Prob}{\mathbb{P}}
\newcommand{\edit}[1]{\emph{\textcolor{blue}{#1}}}
\newcommand*\dif{\mathop{}\!\mathrm{d}}


\renewenvironment{proof}{{\noindent \textbf{\textit{Proof.}}}}
{\hfill $\Box$ \vspace*{0.1in}}
\newenvironment{proofc}{{\noindent \textit{Proof of Claim.}}}
{\hfill $\Box$}



\begin{document}

\begin{LARGE}\begin{center}\bf 16.3 The Fundamental Theorem for Line Integrals\end{center}
\end{LARGE}

\vs\vs



\section{Warm Up}
\textbf{Problem 1:} Find $\int_0^1 f'(x) \dif{x}$ Hint: FToC


\section{Motivation}
We have now unpacked the definition of a line integral and looked at how we can send our curve, which our surface is bounded by, and send it through a vector field! This calculates the amount of work done as a particle travels through down our curve. Today we specifically look at what happens when our vector field is a conservative vector field, meaning it is the gradient of some surface.



\section{Definitions \& Theorems}

\begin{enumerate}
\item The punch line for conservative vector fields is that the line integrals are \textbf{Independent Of Path}. This means that it doesn't matter if my curve twists around like crazy or just goes straight like a line segment. The work done to go from point A to point B is all the same when the vector field is conservative.


\item We therefore get an analogue to the fundamental theorem of calculus because all that matters is our endpoints. Just to remember, the FToC states,

\begin{equation*}
    \int_a^b f'(x) dx = f(b)- f(a).
\end{equation*}

\item Now if our gradient is $\nabla f = \langle f_x, f_y,f_z \rangle $ then $f$ acts as the ``anti-derivative" of the gradient. If $F$ is a conservative vector field then $F = \nabla f$ for some function $f$. We then have the following.

\begin{equation*}
    \int_C F \cdot \dif{r} = \int_C \nabla f \cdot \dif{r} = f(b) - f(a).
\end{equation*}

This is called the fundamental theorem for line integrals. (ONLY FOR CONSERVATIVE VECTOR FIELDS!)

    \item So it all boils down to solving for $f$. We do this by integrating $\int f_x \dif{x}$. Instead of adding a $+C$ for the integrals, since I am integrating with respect to $d$ $x$ I need to add a $+g(y)$. Next take the derivative of this $f$ and set equal to the gradient! (Really just set equal to the missing component (in this case $g'(y)$). Lastly integrate with respect to $y$ to get a $+C$

\item Note if our curve is a closed loop then we know the work done is zero by this theorem! This will be useful for Green's Theorem.



\end{enumerate}




\section{Examples}

\begin{problem}
    Find the work done on $F(x,y) = (2y+1)\hat{i} + (2x+3) \hat{j}$ from $A(0,0)$ to $B(-1,1)$.
\end{problem}

\vspace{9.5cm}

\begin{problem}
    Find the work done on $F(x,y) = xe^{2y}\hat{i} + x^2e^{2y} \hat{j}$ from $A(0,0)$ to $B(-1,1)$.
\end{problem}


\vspace{9.5cm}


\begin{problem}
        Find the work done on $F(x,y) = (x^2+\frac{y}{x})\hat{i} + (y^2+\ln{x}) \hat{j}$ from $A(1,0)$ to $B(e,1)$.
\end{problem}

\vspace{9.5 cm}


\begin{problem}
    Find the work done on $F(x,y) = (x+\tan^{-1}(y))\hat{i} + \left(\frac{x+y}{1+y^2}\right) \hat{j}$ from $A$ to $B$.
\end{problem}

\vspace{9.5 cm}


\begin{problem}
    Find the work done on $F(x,y,z) = yz^2 \hat{i} + xz^2 \hat{j}+ 2xyz \hat{k}$ from $(0,0,1)$ to $(1,3,2)$.
\end{problem}

\newpage

\section{Scratch Work}

\end{document} 